The demand for powerful instruction following and reasoning capability of large language models (LLMs) has promoted rapid development of retrieval-augmented generation (RAG).
The RAG system assists LLM generation by retrieving chunks of query-fit supplementary knowledge from an external database.
%In recent years, large-scale language models (LLMs) have promoted the rapid development of retrieval-augmented generation (RAG) in question-answering and knowledge-based tasks. 
Conventional RAG systems, however, suffer from information insufficiency due to two factors, which are intent-agnostic retrieval and information fragmentation. 
%information fragmentation due to 
%Nevertheless, conventional RAGs suffer from information fragmentation due to blind document chunking.
%as the results of chunking retrieved knowledge chunking.
%However, mainstream RAG systems generally rely on models with huge number of parameters and highly redundant vector retrieval pipelines, resulting in high inference costs, large response delays, and poor evidence controllability. 
%Especially in edge or resource-constrained scenarios, the problems of traditional RAG such as chunk truncation, loss of context, and incomplete knowledge coverage are particularly prominent. 
Our work proposes a RAG framework, termed InSemRAG, that addresses these challenges via an iterative retrieve-and-check mechanism with two supporting modules, an intention-aware retriever (IAR) and semantics-preserving chunking (SPC).
%Specifically, we design a multi-round retrieval mechanism with two supporting modules, intention-aware retriever and semantic-preserving chunk.
IAR implements a dynamic hybrid retrieval method that adaptively weights the retrieval channels based on the query intent, while SPC performs detection and reparation to the damaged evidence chunks to preserve the semantic integrity.
%intention-aware retrieval
%semantic integrity-preserving chunking
%multi-round entropy-based reweighting
%to iteratively optimize the linear combination of
%sparse and dense scores
%implement dynamic hybrid retrieval that 
%To solve these challenges, we propose an RAG framework based on a small-sized LM (SLM) to alleviate the computational latency issue.
%We design the SLM to perform an adaptive retrieval mechanism conditioned on the user queries' needs.
%Specifically, we not only perform dense-sparse query enrichment, but also control the dual retrieval channel through dynamic weighting.
%To preserve the semantic integrity of the retrieved knowledge evidence, we perform semantic damage detection and chunk repair through neighborhood evidence stitching.
%We perform both the retrieval and repair processes iteratively to ensure that we only pass semantically complete evidence to the LLM.
To alleviate the computational latency brought by our iterative mechanism, we leverage small language models (SLMs).
Extensive experiments across several benchmark datasets consistently demonstrate the competitiveness of our method against recent state-of-the-art RAG mechanisms.
Particularly, our method achieves significant gains on multi-hop and evidence-sensitive tasks, with a 2.65-point improvement in F1 on HotPotQA and a 1.5-point increase in accuracy on FEVER.
Our method also achieves competitive performance to Multi-Hop RAG with 4.32$\times$ lower latency with the utilization of SLM.
%Particularly, our method enjoys significant gains on multi-hop and evidence-sensitive tasks (F1=2.65$\uparrow$ on HotPotQA and \% accuracy=1.5 $\uparrow$ on FEVER).
%reduces latency by 90\% compared to LLM-RAG systems while improving evidence quality, interpretability, and response latency by 45\%.
%Our work present a
%To this end, this paper proposes a retrieval-augmented generation framework based on Small Language Model (SLM), SLM-RAG. The framework takes the lightweight model as the core, combined with Hypothetical Query Reasoning, BM25 + Dense Embedding, and forensic evidence completion and coverage physical examination mechanisms, which significantly improves the response efficiency and evidence integrity of the system while maintaining semantic retrieval capabilities. Specifically, SLM is responsible for the semantic-driven tasks of the whole process: generating multi-view hypothetical subqueries during the query phase; assist in block integrity detection and local context completion in the retrieval stage; Evidence screening and answer integration are completed in the rearrangement and generation stages, so as to realize a knowledge enhancement mode of "lightweight autonomy and closed loop of reasoning". Experimental results show that SLM-RAG significantly improves retrieval recall and answer fidelity on multiple public and self-built datasets: while maintaining 1/10 of the computational overhead, the response delay is reduced by about 45\% and the evidence coverage is increased by about 17\% compared with the mainstream LLM-RAG framework. Further analysis shows that SLM-RAG's multi-level reasoning and coverage physical examination mechanism effectively alleviate the problems of text fragmentation, information omission, and hallucination generation. This paper verifies the potential advantages of small language models in complex semantic retrieval and generation tasks, and provides a scalable and low-cost new paradigm for intelligent question answering, enterprise knowledge base retrieval, and personalized agent systems in resource-constrained scenarios.



\begin{figure*}[!htb]
    \centering
    \includegraphics[width=.845\linewidth]{figures/motivation.png}
    \caption{Limitations on the conventional RAG system that relies on a fixed retrieval channel and ignores semantically fragmented evidence chunks. We address these limitations in InSemRAG.}
    \label{fig:motivation}
\end{figure*}